\subsection{An exclusion from closed global one-forms}
\paragraph{A2W (Proved under the stated differential hypothesis).}
Let $X/k$ be smooth, projective and integral of dimension $n\geq2$,
where $k$ is algebraically closed of characteristic $p>0$.
Suppose every global regular one-form on $X$ is closed:
\[
 d:H^0(X,\Omega_X^1)\longrightarrow H^0(X,\Omega_X^2)
 \quad\hbox{is zero}.
\]
If $L$ is very ample, then
\[
 H^0(X,\Omega_X^1\otimes L^{-1})=0.                 \tag{A2W.1}
\]
Consequently, if $n\geq3$, restriction on geometric Picard
groups is injective for every smooth divisor in $|L|$, using
the established reduction of its kernel to a finite $p$-group.
This is a conditional theorem: the closedness hypothesis is
not asserted for all smooth projective varieties in positive
characteristic.

To prove (A2W.1), suppose that
$0\ne\zeta\in H^0(X,\Omega_X^1\otimes L^{-1})$.
Choose a nonzero $s_0\in H^0(X,L)$ and put $\eta=s_0\zeta$.
For every $s\in H^0(X,L)$, both $s\zeta$ and $\eta$ are
global one-forms and hence closed. On the open set $s_0\ne0$
we therefore have
\[
 0=d\bigl((s/s_0)\eta\bigr)=d(s/s_0)\wedge\eta.
\]
Because $L$ gives a closed embedding, its affine coordinate
ratios $s/s_0$ generate the function field $K=k(X)$ over $k$.
Their differentials consequently span the $n$-dimensional
$K$-vector space $\Omega_{K/k}^1$. For a nonzero one-form
$\eta_K$ on a vector space of dimension at least two, the
linear map $v\mapsto v\wedge\eta_K$ is not zero. The
displayed equality for a spanning set is a contradiction.

For the Picard consequence, a nonzero finite $p$-group kernel
would contain a line bundle $M$ of order $p$. The nonzero
logarithmic form $d\log(M)$ restricts to zero tangentially on
the smooth divisor $Y$. The conormal exact sequence and
$H^0(Y,\mathcal O_Y(-Y))=0$ then make its full restriction
zero, putting it in the group (A2W.1). This proves injectivity.

There is a quantitative necessary condition even without the
closedness hypothesis. If a nonzero closed global form $\eta$
vanishes along a very ample divisor $Y$, write $\eta=s_0\zeta$
as above. The differentials of the global forms $s\zeta$ are
$d(s/s_0)\wedge\eta$. Their values over $K$ span the image
of the wedge map by $\eta$, which has rank $n-1$. It follows
that
\[
 \dim_k\operatorname{im}\bigl(
  d:H^0(X,\Omega_X^1)\to H^0(X,\Omega_X^2)\bigr)
  \geq n-1.                                         \tag{A2W.2}
\]
Thus any hypothetical point-kernel example in dimension
$n\geq3$ must have at least $n-1$ independent nonzero
differentials of global one-forms. In particular it must
have a nonzero Hodge-to-de Rham differential $d_1^{1,0}$.
The argument uses only closedness of the original logarithmic
form for this last necessary condition.
